\documentclass[conference]{IEEEtran}

\usepackage[utf8]{inputenc}
\usepackage[T1]{fontenc}
\usepackage[brazil]{babel}
\usepackage{graphicx}
\usepackage{booktabs}
\usepackage{cite}
\usepackage{url}

\title{AnimalSight: Avaliação de Regiões de Interesse\\
na Classificação de Raças de Cães}

\author{
\IEEEauthorblockN{João Meyer Muhlmann}
\IEEEauthorblockA{Universidade Federal do Paraná\\
Trabalho Final de Visão Computacional\\
Novembro de 2026}
}

\begin{document}
\maketitle

\begin{abstract}
Este trabalho investiga se a seleção clássica de uma região de interesse
(ROI) permite reduzir a resolução de entrada do DINOv2 preservando a
acurácia na classificação de cães. A entrada de 112\,$\times$\,112 pixels
reduziu o tempo médio de inferência em aproximadamente 2,9 vezes, com
perda de um ponto percentual de acurácia. Entretanto, a ROI por
gradientes apresentou desempenho inferior ao redimensionamento
convencional em todas as resoluções avaliadas no teste.
\end{abstract}

\begin{IEEEkeywords}
visão computacional, região de interesse, DINOv2, SVM.
\end{IEEEkeywords}

\section{Introdução}
Modelos como o DINOv2 extraem características visuais que podem ser
utilizadas sem ajuste dos seus pesos~\cite{dinov2}. Entretanto, seu custo
de inferência depende da resolução de entrada. O AnimalSight investiga
se um recorte obtido por visão computacional clássica preserva
informações relevantes do animal quando a imagem é reduzida.

\section{Metodologia}
Foram utilizadas cinco raças do Oxford-IIIT Pet~\cite{pets}: beagle,
boxer, chihuahua, pug e samoyed. Mantiveram-se os conjuntos oficiais,
com 500 imagens de treinamento e 499 de teste.

O modelo \texttt{facebook/dinov2-small} permaneceu congelado. O token
CLS da última camada fornece um vetor de 384 características,
padronizado e classificado por um SVM com kernel RBF, $C=10$ e
\texttt{gamma=scale}. Cada representação utiliza um SVM próprio,
treinado nas mesmas 500 imagens, com os parâmetros fixos.

Compararam-se três configurações: \textit{Baseline}, com entrada de
224\,$\times$\,224 pixels; \textit{Reduced}, com entrada de
112\,$\times$\,112; e \textit{ROI Reduced}, que realiza um recorte antes
do processamento em 112\,$\times$\,112. A ROI é uma janela quadrada com
lado igual a 80\% da menor dimensão da imagem, posicionada na região
de maior soma de magnitudes dos gradientes Sobel, suavizados por
filtros gaussianos. Sua seleção não utiliza rótulos ou detectores.

Todas as configurações aplicam redimensionamento, recorte central e
normalização compatíveis com o processador do DINOv2. A resolução dos
tensores efetivamente diminui: de 257 tokens em 224 pixels para 65 em
112, incluindo o CLS. Também foram comparadas entradas de 84, 56,
42 e 28 pixels, mantendo o mesmo algoritmo de ROI.

O benchmark foi executado em CPU Intel Core i7-1065G7, com quatro
threads, lote unitário, duas passagens pelas 499 imagens e cinco
inferências de aquecimento por configuração em cada passagem.
O tempo inclui pré-processamento, ROI quando aplicável, DINOv2 e SVM;
exclui leitura e decodificação dos arquivos, carregamento dos modelos
e treinamento. A medição não utiliza embeddings em cache.

\section{Resultados e Discussão}
A Tabela~\ref{tab:principal} apresenta a comparação principal.
O Reduced alcançou ganho de velocidade de $2{,}90\times$, com acurácia
de 98,80\%, frente a 99,80\% do baseline. A ROI apresentou tempo
próximo, mas reduziu a acurácia para 96,99\%. Os pesos do DINOv2 são
compartilhados e seu tamanho não diminui com a resolução.

\begin{table}[t]
\caption{Comparação principal nas 499 imagens de teste.}
\label{tab:principal}
\centering
\footnotesize
\setlength{\tabcolsep}{4pt}
\begin{tabular}{lrrrr}
\toprule
Configuração & Acc. (\%) & F1 macro & Tempo (ms) & Ganho \\
\midrule
Baseline    & 99,80 & 0,9980 & 202,41 & $1{,}00\times$ \\
Reduced     & 98,80 & 0,9880 & 69,87  & $2{,}90\times$ \\
ROI Reduced & 96,99 & 0,9700 & 67,08  & $3{,}02\times$ \\
\bottomrule
\end{tabular}
\end{table}

\begin{table}[t]
\caption{Acurácia em resoluções menores no conjunto de teste.}
\label{tab:resolucoes}
\centering
\small
\begin{tabular}{crr}
\toprule
Entrada (pixels) & Reduced (\%) & ROI Reduced (\%) \\
\midrule
$112\times112$ & 98,80 & 96,99 \\
$84\times84$   & 95,39 & 94,19 \\
$56\times56$   & 75,15 & 73,55 \\
$42\times42$   & 64,13 & 61,92 \\
$28\times28$   & 57,11 & 54,31 \\
\bottomrule
\end{tabular}
\end{table}

A ROI permaneceu inferior ao Reduced em todas as resoluções
avaliadas (Tabela~\ref{tab:resolucoes}). Ela beneficia algumas imagens:
em 28 pixels, corrigiu 79 erros do Reduced, mas introduziu outros 93.
Uma possível explicação é que gradientes intensos também aparecem
no fundo, enquanto o recorte pode excluir partes importantes do animal.

O estudo é exploratório, limitado a cinco raças e à reutilização do
mesmo conjunto de teste. Não foram medidas novas latências para as
resoluções abaixo de 112 pixels, nem demonstrada significância
estatística das diferenças observadas.

\section{Conclusão e Trabalhos Futuros}
A redução para 112\,$\times$\,112 apresentou uma relação favorável
entre acurácia e tempo de inferência neste experimento. Contudo,
a seleção de ROI por gradientes não preservou melhor a informação
discriminativa que o redimensionamento convencional.

Como trabalho futuro, propõe-se melhorar a delimitação da região do
animal, combinando gradientes com segmentação de primeiro plano e
operações morfológicas. Também podem ser avaliados recortes de
tamanho e formato adaptativos, evitando a exclusão da cabeça ou do
corpo. Essas alternativas devem ser avaliadas considerando tanto a
acurácia quanto o custo adicional do pré-processamento.

\begin{thebibliography}{2}

\bibitem{dinov2}
M. Oquab \textit{et al.}, ``DINOv2: Learning Robust Visual Features
without Supervision,'' arXiv:2304.07193, 2023.
\url{https://arxiv.org/abs/2304.07193}

\bibitem{pets}
O. M. Parkhi, A. Vedaldi, A. Zisserman e C. V. Jawahar,
``Cats and Dogs,'' in \textit{Proc. IEEE Conference on Computer Vision
and Pattern Recognition (CVPR)}, 2012.
\url{https://www.robots.ox.ac.uk/~vgg/data/pets/}

\end{thebibliography}
\end{document}
